% ===================================================================== PREAMBULE COMMUN — Carnet d'entraînement Fiche 9
\documentclass[11pt,a4paper]{article}
\usepackage[utf8]{inputenc}
\usepackage[T1]{fontenc}
\usepackage{lmodern}
\usepackage[french]{babel}
\usepackage{amsmath,amssymb,mathtools}
\usepackage{icomma}
\usepackage{siunitx}
\sisetup{output-decimal-marker={,},per-mode=symbol}
\DeclareSIUnit{\molar}{M}
\DeclareSIUnit{\Molar}{mol\,L^{-1}}
\usepackage[version=4]{mhchem}
\usepackage{xcolor}
\usepackage{colortbl}
\usepackage{tikz}
\usepackage{pgfplots}
\usepackage[most]{tcolorbox}
\usepackage{enumitem}
\usepackage{array}
\usepackage{booktabs}
\usepackage{multicol}
\usepackage{fancyhdr}
\usepackage[hidelinks]{hyperref}
\usetikzlibrary{arrows.meta,positioning,calc,decorations.pathmorphing,
  decorations.markings,angles,quotes,patterns,backgrounds,
  shapes.geometric,shapes.misc,fadings,intersections,fit}
\pgfplotsset{compat=1.18}
\usepackage[a4paper,margin=2.2cm,headheight=26pt,headsep=10pt]{geometry}

% ---------- Palette ----------
\definecolor{primary}{RGB}{150,98,162}
\definecolor{primarydark}{RGB}{92,52,110}
\definecolor{defcol}{RGB}{0,114,178}
\definecolor{thmcol}{RGB}{0,140,120}
\definecolor{formulacol}{RGB}{198,93,9}
\definecolor{remarkcol}{RGB}{120,94,200}
\definecolor{attncol}{RGB}{200,40,55}
\definecolor{excol}{RGB}{0,135,90}
\definecolor{methcol}{RGB}{90,90,100}
\definecolor{cationcol}{RGB}{0,90,190}
\definecolor{anioncol}{RGB}{200,60,55}
\definecolor{cristalcol}{RGB}{110,105,120}
\definecolor{eaucol}{RGB}{120,180,220}
\definecolor{solcol}{RGB}{0,135,90}
\definecolor{kpscol}{RGB}{198,93,9}
\definecolor{qcol}{RGB}{120,94,200}
\definecolor{precipcol}{RGB}{110,70,40}
\definecolor{exercol}{RGB}{150,98,162}
\definecolor{corrcol}{RGB}{0,120,100}

% ---------- Macros ----------
\newcommand{\Kps}{K_{\mathrm{ps}}}
\newcommand{\Ce}[1]{\left[\,\ce{#1}\,\right]}

% ---------- Style des boîtes ----------
\tcbset{boxbase/.style={enhanced,breakable,boxrule=0.9pt,arc=2.5pt,
   left=3mm,right=3mm,top=2.3mm,bottom=2.3mm,
   fonttitle=\bfseries,coltitle=white,toptitle=0.6mm,bottomtitle=0.6mm}}

\newtcolorbox[auto counter]{definition}[1][]{%
  boxbase,colback=defcol!5,colframe=defcol,
  title={\textsc{Définition}~\thetcbcounter\ \ifx\relax#1\relax\relax\else\textmd{--- \itshape #1}\fi}}
\newtcolorbox[auto counter]{theoreme}[1][]{%
  boxbase,colback=thmcol!6,colframe=thmcol,
  title={\textsc{Loi / Propriété}~\thetcbcounter\ \ifx\relax#1\relax\relax\else\textmd{--- \itshape #1}\fi}}
\newtcolorbox[auto counter]{formule}[1][]{%
  boxbase,colback=formulacol!6,colframe=formulacol,
  title={\textsc{Formule}~\thetcbcounter\ \ifx\relax#1\relax\relax\else\textmd{--- \itshape #1}\fi}}
\newtcolorbox{remarque}[1][]{%
  boxbase,colback=remarkcol!6,colframe=remarkcol,
  title={\textsc{Remarque}\ifx\relax#1\relax\relax\else\textmd{ : \itshape #1}\fi}}
\newtcolorbox{attention}[1][]{%
  boxbase,colback=attncol!6,colframe=attncol,
  title={\textsc{Attention}\ifx\relax#1\relax\relax\else\textmd{ : \itshape #1}\fi}}
\newtcolorbox{exemple}[1][]{%
  boxbase,colback=excol!5,colframe=excol,
  title={\textsc{Exemple}\ifx\relax#1\relax\relax\else\textmd{ : \itshape #1}\fi}}
\newtcolorbox{methode}[1][]{%
  boxbase,colback=methcol!7,colframe=methcol,sharp corners,
  title={\textsc{Méthode}\ifx\relax#1\relax\relax\else\textmd{ : \itshape #1}\fi}}
\newtcolorbox{astuce}[1][]{%
  boxbase,colback=primary!7,colframe=primary,
  title={\textsc{Astuce}\ifx\relax#1\relax\relax\else\textmd{ : \itshape #1}\fi}}

\newtcolorbox[auto counter]{exercice}[1][]{%
  boxbase,colback=exercol!4,colframe=exercol,
  title={\textsc{Exercice}~\thetcbcounter\ \ifx\relax#1\relax\relax\else\textmd{--- \itshape #1}\fi}}
\newtcolorbox[auto counter]{correction}[1][]{%
  boxbase,colback=corrcol!4,colframe=corrcol,
  title={\textsc{Corrigé}~\thetcbcounter\ \ifx\relax#1\relax\relax\else\textmd{--- \itshape #1}\fi}}

% ---------- Environnement « où : » ----------
\newenvironment{ou}{%
  \par\smallskip\noindent\textbf{où :}\nopagebreak
  \begin{itemize}[leftmargin=1.8em,topsep=2pt,itemsep=1.5pt,parsep=0pt]}%
  {\end{itemize}}

% ---------- Styles TikZ ----------
\tikzset{
  cat/.style={circle,fill=cationcol,draw=cationcol!50!black,inner sep=0pt,
              minimum size=5mm,font=\bfseries\scriptsize,text=white},
  an/.style={circle,fill=anioncol,draw=anioncol!50!black,inner sep=0pt,
             minimum size=5mm,font=\bfseries\scriptsize,text=white},
  crx/.style={circle,fill=cristalcol,draw=black!50,inner sep=0pt,
              minimum size=5mm,font=\bfseries\scriptsize,text=white},
  ptn/.style={circle,fill=black,inner sep=1.1pt}}

% ---------- En-têtes ----------
\pagestyle{fancy}
\fancyhf{}
\fancyhead[L]{\small\textcolor{primarydark}{\textbf{Fiche 9} — Solubilité et produit de solubilité}}
\fancyhead[R]{\small\textcolor{primarydark}{\textsc{Carnet d'entraînement}}}
\fancyfoot[C]{\small\thepage}
\renewcommand{\headrulewidth}{0.8pt}
\renewcommand{\headrule}{\hbox to\headwidth{\color{primary}\leaders\hrule height \headrulewidth\hfill}}
\usepackage{titlesec}
\titleformat{\section}{\Large\bfseries\color{primarydark}}{\thesection}{0.6em}{}
\titleformat{\subsection}{\large\bfseries\color{primary}}{\thesubsection}{0.6em}{}

% ---------- Bandeau de titre ----------
% #1 = thème/étiquette   #2 = titre principal   #3 = sous-titre
\newcommand{\bandeau}[3]{%
\begin{center}
\begin{tikzpicture}
\node[rectangle,rounded corners=7pt,fill=primary,text=white,
      minimum width=15.6cm,minimum height=1.95cm,align=center,inner sep=8pt]
  {{\large\bfseries #1}\\[3pt]{\Large\bfseries #2}\\[3pt]{\small #3}};
\end{tikzpicture}
\end{center}
\vspace{2mm}}
\begin{document}
\bandeau{Thème 4 — Effet d'ion commun}%
{Corrigé — Niveau Difficile}%
{Réponses détaillées et justifications}

\begin{correction}[récupérer l'argent]
\begin{enumerate}[label=\textbf{\alph*)},leftmargin=2.2em,topsep=2pt,itemsep=2pt]
  \item $s_0=\sqrt{\Kps}=\sqrt{\num{1{,}8e-10}}\approx\qty{1{,}3e-5}{\molar}$. Dans \qty{1{,}0}{\litre} :
        $n_0=\num{1{,}3e-5}\times 1{,}0=\qty{1{,}3e-5}{\mole}$ d'argent dissous.
  \item $s'=\dfrac{\Kps}{[\ce{Cl-}]}=\dfrac{\num{1{,}8e-10}}{0{,}10}=\qty{1{,}8e-9}{\molar}$.
  \item Argent encore dissous : $n'=\num{1{,}8e-9}\times 1{,}0=\qty{1{,}8e-9}{\mole}$, soit une masse
        $m'=n'\,M=\num{1{,}8e-9}\times 108\approx\qty{1{,}9e-7}{\gram}$ (\qty{0{,}19}{\micro\gram}).
  \item Fraction précipitée : $\dfrac{n_0-n'}{n_0}=1-\dfrac{\num{1{,}8e-9}}{\num{1{,}3e-5}}
        \approx 1-\num{1{,}4e-4}=\boxed{99{,}99\,\%}$. Le procédé est \textbf{extrêmement efficace} :
        presque tout l'argent est récupéré sous forme solide (base de la récupération de l'argent
        photographique et de l'argentimétrie).
\end{enumerate}
\end{correction}

\begin{correction}[courbe $s'(C)$ — \ce{CaF2}]
\begin{enumerate}[label=\textbf{\alph*)},leftmargin=2.2em,topsep=2pt,itemsep=2pt]
  \item Ion commun \ce{F-} (2 par unité) : $\Kps=[\ce{Ca^2+}][\ce{F-}]^{2}\approx s'\,C^{2}$, d'où
        $s'=\dfrac{\Kps}{C^{2}}$.
  \item $C=\qty{0{,}050}{\molar}$ : $s'=\dfrac{\num{3{,}9e-11}}{(0{,}050)^{2}}
        =\dfrac{\num{3{,}9e-11}}{\num{2{,}5e-3}}\approx\qty{1{,}6e-8}{\molar}$.\\
        $C=\qty{0{,}20}{\molar}$ : $s'=\dfrac{\num{3{,}9e-11}}{(0{,}20)^{2}}
        =\dfrac{\num{3{,}9e-11}}{\num{4{,}0e-2}}\approx\qty{9{,}8e-10}{\molar}$.
  \item Comme $s'\propto C^{-2}$, multiplier $C$ par 10 \textbf{divise $s'$ par $10^{2}=100$}. Sur le
        graphe log–log, cela se traduit par une droite de \textbf{pente $-2$} (l'exposant de $C$).
\end{enumerate}
\end{correction}

\begin{correction}[ion commun cationique — \ce{PbBr2}]
\begin{enumerate}[label=\textbf{\alph*)},leftmargin=2.2em,topsep=2pt,itemsep=2pt]
  \item La formule $s'=\Kps/C^{\,n}$ suppose que l'\textbf{anion} (ici \ce{Br-}) est l'ion commun apporté
        en excès. Or ici l'ion commun est le \textbf{cation \ce{Pb^2+}} : ce sont les \ce{Br-} qui
        proviennent \emph{entièrement} de la dissolution. On ne peut donc pas appliquer la formule
        telle quelle.
  \item $[\ce{Br-}]=2s'$ (2 bromures par \ce{PbBr2} dissous) ; $[\ce{Pb^2+}]\approx C$. Donc
        \[ \Kps=[\ce{Pb^2+}][\ce{Br-}]^{2}=C\,(2s')^{2}=4\,C\,s'^{2}. \]
  \item $s'^{2}=\dfrac{\Kps}{4C}\Rightarrow \boxed{\ s'=\dfrac{1}{2}\sqrt{\dfrac{\Kps}{C}}\ }$.
  \item $s'=\dfrac{1}{2}\sqrt{\dfrac{\num{9{,}2e-6}}{0{,}10}}=\dfrac{1}{2}\sqrt{\num{9{,}2e-5}}
        =\dfrac{1}{2}\times\num{9{,}6e-3}\approx\boxed{\qty{4{,}8e-3}{\molar}}$.
        Comparé à $s_0=\qty{1{,}3e-2}{\molar}$ (eau pure), la solubilité est divisée par $\sim 2{,}7$ :
        l'effet d'ion commun agit toujours dans le sens d'une \emph{baisse}, mais plus modérément ici
        (dépendance en $\sqrt{1/C}$, plus douce que le $1/C^{2}$ précédent).
\end{enumerate}
\end{correction}

\end{document}
